\section{Теоретическое обоснование и проектирование решения}

\subsection{Принципы и логическая модель решения}

\subsubsection{Модель многопоточных вычислений: работа, глубина,
критический путь}

Для теоретического обоснования эффективности механизма work-stealing
необходима формальная модель, позволяющая описывать параллельные
вычисления и рассуждать об их производительности независимо от
конкретной реализации. В данном разделе излагается модель,
предложенная Блюмофе и Лейзерсоном~\cite{Blumofe1999} и ставшая
стандартной теоретической основой для анализа параллельных
планировщиков задач.

В модели Блюмофе и Лейзерсона параллельное вычисление представляется
в виде ориентированного ациклического графа $G = (V, E)$, где
множество вершин $V$ соответствует элементарным вычислительным
шагам (инструкциям или задачам), а множество рёбер $E$ задаёт
отношения зависимости между ними: ребро $(u, v) \in E$ означает,
что задача $v$ не может начать выполнение до завершения задачи $u$.

Граф вычислений обладает единственной начальной вершиной (истоком),
соответствующей точке входа в параллельную программу, и единственной
конечной вершиной (стоком), соответствующей точке её завершения.
Вершины, не связанные отношением зависимости, могут выполняться
параллельно на различных процессорах. Таким образом, граф $G$
полностью описывает структуру параллелизма вычисления: он определяет,
какие задачи могут выполняться одновременно, а какие должны
выполняться последовательно.

Рассмотрим в качестве иллюстрации алгоритм параллельного вычисления
суммы массива методом «разделяй и властвуй». Массив рекурсивно
делится пополам, левая и правая половины суммируются параллельно,
после чего результаты складываются. Соответствующий граф вычислений
имеет древовидную структуру: каждый внутренний узел соответствует
операции сложения двух частичных сумм, а листья --- суммированию
отдельных элементов. Рёбра графа направлены от листьев к корню,
отражая зависимость каждой операции сложения от результатов
нижележащих задач.

На основе графа $G$ вводятся два ключевых параметра, определяющих
теоретические пределы производительности параллельного
вычисления~\cite{Blumofe1999}.

Работа $T_1$ --- суммарное время выполнения всех вершин графа
при использовании одного процессора, то есть длина оптимального
последовательного выполнения:

\begin{equation}
    T_1 = \sum_{v \in V} w(v),
    \label{eq:work}
\end{equation}

где $w(v)$ --- время выполнения задачи $v$. Работа представляет
собой нижнюю оценку суммарного объёма вычислений, который необходимо
выполнить независимо от числа процессоров.

Глубина (критический путь) $T_\infty$ --- длина наидлиннейшего
пути в графе $G$ от истока до стока при неограниченном числе
процессоров:

\begin{equation}
    T_\infty = \max_{p \in \text{paths}(G)} \sum_{v \in p} w(v),
    \label{eq:span}
\end{equation}

где максимум берётся по всем путям $p$ от истока до стока.
Глубина определяет нижнюю оценку времени выполнения вычисления,
которую невозможно преодолеть никаким числом процессоров: задачи,
лежащие на критическом пути, вынуждены выполняться последовательно
в силу отношений зависимости.

Параметры $T_1$ и $T_\infty$ задают два нижних предела для времени
выполнения $T_P$ вычисления на $P$ процессорах~\cite{Blumofe1999}:

\begin{equation}
    T_P \geq \frac{T_1}{P}, \quad T_P \geq T_\infty.
    \label{eq:lower-bounds}
\end{equation}

Первое неравенство отражает тот факт, что $P$ процессоров не могут
выполнить работу $T_1$ быстрее, чем за $T_1 / P$. Второе неравенство
утверждает, что никакой планировщик не может выполнить вычисление
быстрее, чем длина критического пути.

Блюмофе и Лейзерсон доказали, что жадный алгоритм work-stealing
обеспечивает время выполнения, близкое к теоретическому
минимуму~\cite{Blumofe1999}. А именно, при выполнении вычисления
с работой $T_1$ и глубиной $T_\infty$ на $P$ процессорах
планировщик work-stealing гарантирует:

\begin{equation}
    T_P \leq \frac{T_1}{P} + O(T_\infty).
    \label{eq:work-stealing-bound}
\end{equation}

Данная оценка является оптимальной с точностью до константного
множителя: первое слагаемое соответствует нижней оценке,
обусловленной объёмом работы, второе --- нижней оценке,
обусловленной глубиной вычисления. Важно отметить, что данная
гарантия обеспечивается без какой-либо информации о структуре
графа вычислений: планировщик принимает решения о распределении
задач динамически, исключительно на основе текущего состояния
очередей потоков.

Для количественной оценки эффективности параллельного выполнения
вводится понятие ускорения (speedup) --- отношения времени
последовательного выполнения к времени параллельного выполнения
на $P$ процессорах~\cite{Herlihy2020}:

\begin{equation}
    S_P = \frac{T_1}{T_P}.
    \label{eq:speedup}
\end{equation}

Идеальным (линейным) ускорением называется случай $S_P = P$,
когда использование $P$ процессоров даёт ровно $P$-кратное
ускорение. На практике достичь линейного ускорения невозможно
из-за наличия последовательных участков вычисления, накладных
расходов на синхронизацию и эффектов иерархии памяти.

Фундаментальное ограничение на достижимое ускорение устанавливает
закон Амдала~\cite{Amdahl1967}. Пусть доля последовательной части
вычисления (то есть части, которая не может быть распараллелена)
составляет $\alpha \in [0, 1]$. Тогда максимальное ускорение
при использовании $P$ процессоров ограничено сверху:

\begin{equation}
    S_P \leq \frac{1}{\alpha + \dfrac{1 - \alpha}{P}}.
    \label{eq:amdahl}
\end{equation}

При $P \to \infty$ данное выражение стремится к $1 / \alpha$,
то есть максимальное ускорение ограничено величиной, обратной
доле последовательной части. Так, если 10\% вычисления является
последовательным ($\alpha = 0{,}1$), то максимально достижимое
ускорение не превышает 10 независимо от числа процессоров.

Закон Амдала подчёркивает значимость минимизации последовательных
участков в параллельных программах. Именно поэтому накладные
расходы на синхронизацию, неизбежно возникающие при использовании
разделяемых структур данных, столь критичны для производительности:
каждый момент ожидания блокировки фактически увеличивает
последовательную долю вычисления и тем самым ограничивает
достижимое ускорение.

В контексте разрабатываемой библиотеки закон Амдала определяет
ключевое требование к реализации: накладные расходы на управление
очередями задач и механизм кражи должны быть минимизированы,
чтобы не увеличивать эффективную долю последовательной части
вычисления. Именно это требование обусловливает использование
атомарных операций вместо блокирующей синхронизации в реализации
двусторонней очереди.